% --- Archon physics formalization source begin ---
% archon:physics
% archon:chemistry
% archon:covers IChO2026Problems/problem_icho_2026_t2_a3.lean
% archon:source-report reports/icho_2026/problem_icho_2026_t2_a3.source.json
% archon:problem-id icho_2026_t2
% archon:part-id T2-A3
% archon:previous-part icho_2026_t2_a2
% archon:previous-part-policy natural_language_prerequisite_only

\section{Icho Chemistry problem icho\_2026\_t2\_a3}

\paragraph{Problem source.}
\#\# Source
58th International Chemistry Olympiad, Tashkent, Uzbekistan, 2026.
Official-materials index: https://www.icho2026.uz/problems
English problem PDF: ../icho\_2026\_source/raw/theory\_problem.pdf
English solution/rubric PDF: ../icho\_2026\_source/raw/theory\_solution.pdf
Problem PDF page 17; printed local question page 3; solution starts on PDF page 17.
The page PNG is primary visual evidence for formulas, structures, tables, and figures.

\#\# T2. Kinetics of the Belousov-Zhabotinsky Reaction
T2. Kinetics of the Belousov-Zhabotinsky reaction 6\% When the chemicals KBrO3, CH2(COOH)2 (malonic acid, – MA), Ce(SO4)2, and H2SO4 are mixed, the concentration, 𝐶, of Ce species and the colour of the solution oscillates between yellow (due to Ce4+) and colourless (due to Ce3+). A schematic picture is given below: 2.1 Write the overall balanced reaction equation of the Belousov-Zhabotinsky reaction (BZ reaction). Note: in this reaction, MA is oxidised to CO2, BrO− 3 is reduced to Br−, and Ce(IV) is a catalyst. The mechanism of the BZ reaction is shown below (BMA −CHBr(COOH)2): Process A ⎧ ⎪ ⎪ ⎨ ⎪ ⎪ ⎩ (1) HBrO2 + BrO− 3 + H+ 𝑘1 −−→2BrO⋅ 2 + H2O (2) BrO⋅ 2 + Ce3+ + H+ 𝑘2 −−→HBrO2 + Ce4+ (3) 2HBrO2 𝑘3 −−→BrO− 3 + HBrO + H+ 𝑘1 = 1.0 × 104 M−2 s−1 𝑘2 = 6.2 × 104 M−2 s−1 𝑘3 = 4.0 × 107 M−1 s−1 Process B ⎧ ⎪ ⎪ ⎨ ⎪ ⎪ ⎩ (4) HBrO2 + Br−+ H+ 𝑘4 −−→2HBrO (5) BrO− 3 + Br−+ 2H+ 𝑘5 −−→HBrO + HBrO2 (6) HBrO + MA 𝑘6 −−→BMA + H2O 𝑘4 = 2.0 × 109 M−2 s−1 𝑘5 = 2.1 M−3 s−1 𝑘6 = 8.2 M−1 s−1 Process C \{(7) Ce4+ + BMA 𝑘7 −−→Ce3+ + Br−+ other products 𝑘7 = 1.0 × 102 M−1 s−1 In the BZ reaction, Process C occurs continuously, while Processes A and B alternate with each other. Assume that: - when Process A occurs the solution is yellow and Process B practically does not occur; - when Process B occurs the solution is colourless and Process A practically does not occur; - [BrO3−]0 = 0.06 M, [MA]0 = 0.1 M, [H+]0 = 0.8 M, [Ce4+]0 = 0.001 M; - the concentrations of the reactants (except Ce4+) and pH are maintained constant throughout the BZ reaction. 2.2 Using the steady-state approximation, calculate the stationary molar concentrations of HBrO2 in Processes A and B, [HBrO2]A and [HBrO2]B. If you could not calculate [HBrO2]A and [HBrO2]B, you can use the values [HBrO2]A = 1 × 10−5 M and [HBrO2]B = 1 × 10−10 M for the following calculations. During the BZ reaction, the concentrations of HBrO2 and Br−oscillate strictly in relation to each other, and this is sketched in the phase portrait of [HBrO2] vs [Br−]: Here [Br−]max is the maximum concentration and [Br−]critical is the critical concentration of Br−at which there is a switch over in processes from A to B or from B to A. To switch Process A to Process B, the reaction rate of the elementary step (4) must exceed that of the elementary step (1) and vice versa.

\#\# Current subquestion 2.3
Calculate [Br−]critical.

\#\# Dataset metadata
IChO 2026; T2; raw rubric points=4; kind=theory; formalization\_ready=true.

\paragraph{Shared problem context.}
T2. Kinetics of the Belousov-Zhabotinsky reaction 6\% When the chemicals KBrO3, CH2(COOH)2 (malonic acid, – MA), Ce(SO4)2, and H2SO4 are mixed, the concentration, 𝐶, of Ce species and the colour of the solution oscillates between yellow (due to Ce4+) and colourless (due to Ce3+). A schematic picture is given below: 2.1 Write the overall balanced reaction equation of the Belousov-Zhabotinsky reaction (BZ reaction). Note: in this reaction, MA is oxidised to CO2, BrO− 3 is reduced to Br−, and Ce(IV) is a catalyst. The mechanism of the BZ reaction is shown below (BMA −CHBr(COOH)2): Process A ⎧ ⎪ ⎪ ⎨ ⎪ ⎪ ⎩ (1) HBrO2 + BrO− 3 + H+ 𝑘1 −−→2BrO⋅ 2 + H2O (2) BrO⋅ 2 + Ce3+ + H+ 𝑘2 −−→HBrO2 + Ce4+ (3) 2HBrO2 𝑘3 −−→BrO− 3 + HBrO + H+ 𝑘1 = 1.0 × 104 M−2 s−1 𝑘2 = 6.2 × 104 M−2 s−1 𝑘3 = 4.0 × 107 M−1 s−1 Process B ⎧ ⎪ ⎪ ⎨ ⎪ ⎪ ⎩ (4) HBrO2 + Br−+ H+ 𝑘4 −−→2HBrO (5) BrO− 3 + Br−+ 2H+ 𝑘5 −−→HBrO + HBrO2 (6) HBrO + MA 𝑘6 −−→BMA + H2O 𝑘4 = 2.0 × 109 M−2 s−1 𝑘5 = 2.1 M−3 s−1 𝑘6 = 8.2 M−1 s−1 Process C \{(7) Ce4+ + BMA 𝑘7 −−→Ce3+ + Br−+ other products 𝑘7 = 1.0 × 102 M−1 s−1 In the BZ reaction, Process C occurs continuously, while Processes A and B alternate with each other. Assume that: - when Process A occurs the solution is yellow and Process B practically does not occur; - when Process B occurs the solution is colourless and Process A practically does not occur; - [BrO3−]0 = 0.06 M, [MA]0 = 0.1 M, [H+]0 = 0.8 M, [Ce4+]0 = 0.001 M; - the concentrations of the reactants (except Ce4+) and pH are maintained constant throughout the BZ reaction. 2.2 Using the steady-state approximation, calculate the stationary molar concentrations of HBrO2 in Processes A and B, [HBrO2]A and [HBrO2]B. If you could not calculate [HBrO2]A and [HBrO2]B, you can use the values [HBrO2]A = 1 × 10−5 M and [HBrO2]B = 1 × 10−10 M for the following calculations. During the BZ reaction, the concentrations of HBrO2 and Br−oscillate strictly in relation to each other, and this is sketched in the phase portrait of [HBrO2] vs [Br−]: Here [Br−]max is the maximum concentration and [Br−]critical is the critical concentration of Br−at which there is a switch over in processes from A to B or from B to A. To switch Process A to Process B, the reaction rate of the elementary step (4) must exceed that of the elementary step (1) and vice versa.

\paragraph{Current subquestion.}
Calculate [Br−]critical.

\paragraph{Recorded answer/context.}
During the switching process the following equality is observed: 𝑘1[HBrO2][BrO− 3 ][H+] = 𝑘4[HBrO2][Br−][H+] (eq.3.1) [Br−]𝑐𝑟𝑖𝑡𝑖𝑐𝑎𝑙= 𝑘1 𝑘4 [BrO− 3 ] (eq.3.2) [Br−]𝑐𝑟𝑖𝑡𝑖𝑐𝑎𝑙= 1.0×104 2.0×109 ⋅0.06 = 3.0 × 10−7 M [Br−]𝑐𝑟𝑖𝑡𝑖𝑐𝑎𝑙= 3.0 × 10−7 M For fully correct equation 3.1: 1 point. For fully correct equation 3.2: 2 points. For calculated [Br−]critical value: 1 point. For correct calculation without formula derivation: full marks. 4.0 pt

\paragraph{Figure/image paths.}
\begin{itemize}
\item ../icho\_2026\_source/image/T2\_page-3.png
\item ../icho\_2026\_source/image/T2\_page-2.png
\end{itemize}

\paragraph{Reusable previous-part conclusions.}
\begin{itemize}
\item Source T2-A2. Question: Using the steady-state approximation, calculate the stationary molar concentrations of HBrO2 in Processes A and B, [HBrO2]A and [HBrO2]B. Reusable conclusions: 𝑑[BrO⋅ 2] 𝑑𝑡 = 2𝑘1[HBrO2][BrO− 3 ][H+] −𝑘2[BrO⋅ 2][Ce3+][H+] = 0 (eq.2.1) 𝑑[HBrO2] 𝑑𝑡 = −𝑘1[HBrO2][BrO− 3 ][H+]+𝑘2[BrO⋅ 2][Ce3+][H+]−2𝑘3[HBrO2]2 = 0 (eq.2.2) When adding equations 2.1 and 2.2 we get the following: 𝑘1[HBrO2][BrO− 3 ][H+] −2𝑘3[HBrO2]2 = 0 (eq.2.3) [HBrO2]𝐴= 𝑘1 2𝑘3 [BrO− 3 ][H+] (eq.2.4) [HBrO2]𝐴= 1.0×104 2×4.0×107 × 0.06 × 0.8 = 6.0 × 10−6 M The steady-state approximation for Process B: 𝑑[HBrO2] 𝑑𝑡 = −𝑘4[HBrO2][Br−][H+] + 𝑘5[BrO− 3 ][Br−][H+]2 = 0 (eq.2.5) [HBrO2]𝐵= 𝑘5 𝑘4 [BrO− 3 ][H+] (eq.2.6) [HBrO2]𝐵= 2.1 2.0×109 × 0.06 × 0.8 = 5.04 × 10−11 M [HBrO2]𝐴= 6.0 × 10−6 M [HBrO2]𝐵= 5.04 × 10−11 M For fully correct equations 2.1, 2.2 and 2.5: 2 points each. If only one incorrect coefficient in front of the rate constant in equations 2.1, 2.2 and 2.5: 1 point each. For corresponding equations 2.3, 2.4 and 2.6: 1 point each. For calculated [HBrO2]A and [HBrO2]B values: 1 point each. For correct calculation without formula derivation: full marks. 11.0 pt Policy: natural\_language\_prerequisite\_only; the current target and later answers must not be assumed
\end{itemize}

\begin{definition}\leanok
[T2-A3 switching data]
\label{def:chem:icho_2026_t2_a3:data}
\lean{IChO2026Problems.T2A3.BZSwitchingData}
\uses{def:chem:bz:state, def:chem:bz:kinetic_parameters}
The switching datum uses the shared BZ state and kinetic parameters.  Let
$k_1,k_4$, and $b=[\mathrm{BrO_3^-}]$ be readouts with the
source values $k_1=10^4$, $k_4=2\cdot10^9$, and $b=0.06$.  At a switching
point, retain $x=[\mathrm{HBrO_2}]$, $h=[\mathrm{H^+}]$, and
$y=[\mathrm{Br^-}]_{\mathrm{critical}}$ explicitly; assume $x>0$ and
$h>0$ so that the common rate factors may legitimately be cancelled.
\end{definition}
\begin{proof}

\leanok
The record separates the supplied rate constants and bromate concentration
from the switching-state concentrations, so the cancellation conditions remain
visible hypotheses rather than implicit physical assumptions.
\end{proof}

\begin{definition}\leanok
[Process-A elementary rate]
\label{def:chem:icho_2026_t2_a3:process_a_rate}
\lean{IChO2026Problems.T2A3.processARate}
\uses{def:chem:icho_2026_t2_a3:data, def:chem:bz:rate1}
The rate of step (1) is the mass-action product
$k_1[\mathrm{HBrO_2}][\mathrm{BrO_3^-}][\mathrm{H^+}]$ at the switching
state.
\end{definition}
\begin{proof}

\leanok
This is the displayed source rate law evaluated from the switching-data
fields.
\end{proof}

\begin{definition}\leanok
[Process-B elementary rate]
\label{def:chem:icho_2026_t2_a3:process_b_rate}
\lean{IChO2026Problems.T2A3.processBRate}
\uses{def:chem:icho_2026_t2_a3:data, def:chem:bz:rate4}
The rate of step (4) is the mass-action product
$k_4[\mathrm{HBrO_2}][\mathrm{Br^-}][\mathrm{H^+}]$ at the switching state.
\end{definition}
\begin{proof}

\leanok
This is the displayed source rate law evaluated from the same switching-data
fields.
\end{proof}

\begin{theorem}\leanok
[Cancellation at the switching point]
\label{lem:chem:icho_2026_t2_a3:switching_balance}
\lean{IChO2026Problems.T2A3.switching_rate_balance}
\uses{def:chem:icho_2026_t2_a3:data, def:chem:icho_2026_t2_a3:process_a_rate, def:chem:icho_2026_t2_a3:process_b_rate}
If the two elementary rates are equal,
\[
 k_1xbh=k_4xyh,
\]
then $k_1b=k_4y$.
\end{theorem}
\begin{proof}


\leanok
Subtract the two sides and factor the difference as
$xh(k_1b-k_4y)=0$.  Since $x$ and $h$ are positive, their product is
nonzero, so the final factor vanishes.
\end{proof}

% SOURCE: IChO 2026 Theory Solutions, T2.3 (read from references/icho_2026_theory_ready.jsonl)
% SOURCE QUOTE: "[Br−]𝑐𝑟𝑖𝑡𝑖𝑐𝑎𝑙= 3.0 × 10−7 M"
\begin{theorem}\leanok
[Critical bromide concentration]
\label{thm:physics:icho_2026_t2_a3:target}
\lean{IChO2026Problems.T2A3.critical_bromide_concentration}
\uses{def:chem:icho_2026_t2_a3:data, lem:chem:icho_2026_t2_a3:switching_balance}
\textit{Source: IChO 2026 Theory Solutions, T2.3.}
At the switch between Processes A and B,
\[
 [\mathrm{Br^-}]_{\mathrm{critical}}
 =\frac{k_1}{k_4}[\mathrm{BrO_3^-}]
 =3.0\cdot10^{-7}\,\mathrm{M}.
\]
This theorem does not assume either numerical answer from T2-A2.
\end{theorem}
% SOURCE QUOTE PROOF: "During the switching process the following equality is observed: 𝑘1[HBrO2][BrO− 3 ][H+] = 𝑘4[HBrO2][Br−][H+] (eq.3.1)"
\begin{proof}


\leanok
The switching-balance lemma gives $k_1b=k_4y$.  Since the supplied $k_4$ is
positive, divide to obtain $y=(k_1/k_4)b$.  Direct substitution gives
$(10^4/(2\cdot10^9))\cdot0.06=3.0\cdot10^{-7}$.  The argument uses the
current rate equality only, respecting the source policy that T2-A2 is a
natural-language prerequisite rather than a formal assumption.
\end{proof}
% --- Archon physics formalization source end ---
